\lessonlabel{1}
\section{What Is a Database?}

\begin{note}
\textbf{Learning outcomes.} By the end of this lesson you should be able to:
\begin{itemize}
  \item distinguish between data and information;
  \item explain the problems of storing data in ordinary files;
  \item describe what a database management system (DBMS) does for us;
  \item decide when a database is the right tool for a problem.
\end{itemize}
\end{note}

\subsection{Data and information}

\term{Data} are raw facts: numbers, words, dates, or images that have been recorded but not yet interpreted. On their own, the value \texttt{23} or the text \texttt{Okafor} says very little.

\term{Information} is data that has been processed, organised, or placed in a context so that it answers a question. The statement ``Student 1043 scored 23 out of 30 in Database Systems'' is information, because it describes a specific person and a specific result.

The step from data to information is what databases are built to support. A database stores data in a structured way so that it can be retrieved, combined, and checked to produce information quickly and reliably.

\subsection{Storing data in files}

Before databases, most data was kept in ordinary files, one file per purpose. A registry might keep a file of students, another of courses, and another of results. Each program that needed the data read its own file.

This approach works for small amounts of data, but it creates four recurring problems:

\begin{enumerate}
  \item \term{Redundancy.} The same fact is copied into several files. A student's address may appear in the student file, the fees file and the hostel file.
  \item \term{Inconsistency.} When one copy is updated and the others are not, the files disagree. The fees file says one address, the hostel file says another.
  \item \term{Difficult access.} Answering a new question, such as ``Which students owe fees and live in hostel B?'', usually means writing a new program to read and combine the files.
  \item \term{Concurrent use.} If two people update the same file at the same time, one update can be lost.
\end{enumerate}

\begin{example}
A school keeps a \texttt{students.txt} file and a \texttt{fees.txt} file. A student changes address. The office updates \texttt{students.txt} but forgets \texttt{fees.txt}. The next fee reminder is posted to the old address, and nobody notices the mismatch until a complaint arrives.
\end{example}

\subsection{What a database is}

A \term{database} is an organised collection of related data, stored so that many users and programs can share it. A \term{database management system} (DBMS) is the software that creates, stores, protects and queries the database. Examples include MySQL, PostgreSQL, Oracle and SQLite.

A DBMS provides several services that file-based systems do not:
\begin{itemize}
  \item \textbf{Single copy of data.} Each fact is stored once, so updates happen in one place.
  \item \textbf{Querying.} Questions are written in a query language, usually SQL, and the DBMS finds the answer.
  \item \textbf{Integrity.} Rules such as ``every result must belong to an existing student'' are enforced automatically.
  \item \textbf{Concurrency and recovery.} Many users can work at once, and the database can recover from a crash without losing committed work.
  \item \textbf{Security.} Access can be limited to particular users and particular actions.
\end{itemize}

\subsection{When a database is the right tool}

A database is worth its cost when the data is shared, changes often, and must stay correct. A single list of ten names that one person edits once a month does not need one. A university's student records, which many offices read and update every day, do.

\begin{keytest}
\textbf{Key test.} Ask two questions. Will more than one program or person need this data? Must the data remain correct when it changes? If both answers are yes, a database is almost certainly the better choice.
\end{keytest}

\subsection{Exercises}

\begin{enumerate}
  \item Give one example each of data and information from your own course registration.
  \item A bookshop keeps stock in a spreadsheet and sales in a separate notebook. Name two problems this may cause.
  \item Explain why a DBMS stores each fact once, and why this helps when a fact changes.
\end{enumerate}

\subsection{Key terms}

\begin{tabular}{@{}p{3.6cm}p{10.4cm}@{}}
\toprule
\textbf{Term} & \textbf{Meaning} \\
\midrule
Data & Raw recorded facts without interpretation. \\
Information & Data organised to answer a question. \\
DBMS & Software that stores, protects and queries a database. \\
Redundancy & The same fact stored in more than one place. \\
Integrity & The rules that keep data correct. \\
\bottomrule
\end{tabular}
